\chapter*{Conclusion Générale}
\addcontentsline{toc}{chapter}{Conclusion Générale}
\markboth{Conclusion générale}{}

Le présent travail de fin d'études a permis de concevoir, structurer et réaliser un prototype opérationnel de niveau industriel de plateforme mobile de covoiturage inter-wilayas sécurisée, baptisée \textit{RohWinBghit}. Au terme de ce cycle de développement et de recherche, nous présentons le bilan synthétique du projet, la confrontation de nos hypothèses de départ aux résultats expérimentaux, ainsi que les limites identifiées et les perspectives d'évolution.

\section*{Bilan du projet}
\addcontentsline{toc}{section}{Bilan du projet}

Au terme de ce cycle d'ingénierie logicielle, les objectifs SMART fixés au démarrage du projet ont été atteints. La contribution de ce travail s'organise autour de trois axes de valeur~:
\begin{enumerate}[leftmargin=1.8em]
  \item \textbf{Contribution d'ingénierie logicielle~:} La structuration d'une architecture multi-services découplée, conteneurisée et robuste, tolérante aux pannes réseau routières par le biais d'un cache mobile intelligent et de mécanismes de résilience réseau.
  \item \textbf{Contribution de sécurité biométrique~:} L'intégration au sein de l'écosystème d'identité algérien d'un pipeline complet d'authentification biométrique d'anti-spoofing à base d'apprentissage profond, éliminant les risques de vols d'identité.
  \item \textbf{Contribution économique (Arrêté 1275)~:} L'élaboration d'un modèle économique d'intermédiation financière viable intégrant le paiement interbancaire local et fondé sur un plan d'affaires rigoureux, ouvrant la voie à la labellisation en startup.
\end{enumerate}

\section*{Confrontation et validation des hypothèses}
\addcontentsline{toc}{section}{Confrontation et validation des hypothèses}

\begin{itemize}[leftmargin=1.8em]
  \item \textbf{Hypothèse H1 (Pipeline biométrique KYC vs Documentaire)~:} \textit{Partiellement validée.} Bien que l'architecture biométrique à trois niveaux soit pleinement opérationnelle et se montre sémantiquement plus robuste qu'un contrôle documentaire statique passif face aux scénarios d'attaque de rejeu, l'absence d'évaluation formelle et exhaustive sur des bases de données de référence normalisées ne permet pas d'établir sa supériorité quantitative exacte. Cette évaluation expérimentale constitue un axe de recherche prioritaire.
  \item \textbf{Hypothèse H2 (Paiement local et Utilisabilité SUS)~:} \textit{Validée.} L'obtention d'un score moyen normalisé de 71,6 sur 100 à l'évaluation SUS, intégrant les parcours d'authentification trilingues et de simulation de la monétique nationale (CIB/Edahabia), valide l'acceptabilité ergonomique de la solution par les usagers algériens.
  \item \textbf{Hypothèse H3 (Tolérance aux déconnexions et Disponibilité API)~:} \textit{Validée.} Les simulations de charge sous contraintes de connectivité intermittente démontrent que l'utilisation conjointe d'idempotence applicative, de caches persistants et de files de retentatives asynchrones permet de maintenir un taux d'exécution des transactions API de 99,9\,\% pour les flux transactionnels clés.
\end{itemize}

\section*{Limites et perspectives d'évolution}
\addcontentsline{toc}{section}{Limites et perspectives d'évolution}

\subsection*{Limites identifiées}
\begin{enumerate}[leftmargin=1.8em]
  \item La dépendance de l'écosystème de paiement envers l'environnement d'homologation de la SATIM, imposant de maintenir le projet en sandbox monétique dans l'attente de l'homologation formelle de production.
  \item Le biais méthodologique d'échantillonnage de l'évaluation SUS, restreint à une population d'utilisateurs hautement technophiles en milieu universitaire.
  \item L'absence d'audit de sécurité applicatif complet réalisé par un tiers externe certifié en cybersécurité.
\end{enumerate}

\subsection*{Perspectives de développement}
\begin{itemize}[leftmargin=1.8em]
  \item \textbf{À court terme~:} Déploiement d'un projet pilote d'évaluation opérationnelle restreint à l'axe routier Tlemcen-Oran, stabilisation de l'isolation transactionnelle de la suite de tests d'intégration, et obtention du label ministériel startup sous l'arrêté 1275.
  \item \textbf{À moyen terme~:} Intégration d'un module d'IA prédictive d'aide à la décision pour le calcul en temps réel de la suggestion tarifaire par wilaya, et négociation de partenariats logistiques de type B2B avec des institutions universitaires et industrielles.
  \item \textbf{À long terme~:} Lancement d'une phase d'expansion commerciale de la plateforme vers les pays d'Afrique du Nord et du Moyen-Orient (région MENA) soumis à des contraintes de mobilité similaires.
\end{itemize}
